% Options for packages loaded elsewhere
% Options for packages loaded elsewhere
\PassOptionsToPackage{unicode}{hyperref}
\PassOptionsToPackage{hyphens}{url}
\PassOptionsToPackage{dvipsnames,svgnames,x11names}{xcolor}
%
\documentclass[
  11pt,
]{article}
\usepackage{xcolor}
\usepackage[margin=1in]{geometry}
\usepackage{amsmath,amssymb}
\setcounter{secnumdepth}{-\maxdimen} % remove section numbering
\usepackage{iftex}
\ifPDFTeX
  \usepackage[T1]{fontenc}
  \usepackage[utf8]{inputenc}
  \usepackage{textcomp} % provide euro and other symbols
\else % if luatex or xetex
  \usepackage{unicode-math} % this also loads fontspec
  \defaultfontfeatures{Scale=MatchLowercase}
  \defaultfontfeatures[\rmfamily]{Ligatures=TeX,Scale=1}
\fi
\usepackage{lmodern}
\ifPDFTeX\else
  % xetex/luatex font selection
\fi
% Use upquote if available, for straight quotes in verbatim environments
\IfFileExists{upquote.sty}{\usepackage{upquote}}{}
\IfFileExists{microtype.sty}{% use microtype if available
  \usepackage[]{microtype}
  \UseMicrotypeSet[protrusion]{basicmath} % disable protrusion for tt fonts
}{}
\makeatletter
\@ifundefined{KOMAClassName}{% if non-KOMA class
  \IfFileExists{parskip.sty}{%
    \usepackage{parskip}
  }{% else
    \setlength{\parindent}{0pt}
    \setlength{\parskip}{6pt plus 2pt minus 1pt}}
}{% if KOMA class
  \KOMAoptions{parskip=half}}
\makeatother
% Make \paragraph and \subparagraph free-standing
\makeatletter
\ifx\paragraph\undefined\else
  \let\oldparagraph\paragraph
  \renewcommand{\paragraph}{
    \@ifstar
      \xxxParagraphStar
      \xxxParagraphNoStar
  }
  \newcommand{\xxxParagraphStar}[1]{\oldparagraph*{#1}\mbox{}}
  \newcommand{\xxxParagraphNoStar}[1]{\oldparagraph{#1}\mbox{}}
\fi
\ifx\subparagraph\undefined\else
  \let\oldsubparagraph\subparagraph
  \renewcommand{\subparagraph}{
    \@ifstar
      \xxxSubParagraphStar
      \xxxSubParagraphNoStar
  }
  \newcommand{\xxxSubParagraphStar}[1]{\oldsubparagraph*{#1}\mbox{}}
  \newcommand{\xxxSubParagraphNoStar}[1]{\oldsubparagraph{#1}\mbox{}}
\fi
\makeatother


\usepackage{longtable,booktabs,array}
\newcounter{none} % for unnumbered tables
\usepackage{calc} % for calculating minipage widths
% Correct order of tables after \paragraph or \subparagraph
\usepackage{etoolbox}
\makeatletter
\patchcmd\longtable{\par}{\if@noskipsec\mbox{}\fi\par}{}{}
\makeatother
% Allow footnotes in longtable head/foot
\IfFileExists{footnotehyper.sty}{\usepackage{footnotehyper}}{\usepackage{footnote}}
\makesavenoteenv{longtable}
\usepackage{graphicx}
\makeatletter
\newsavebox\pandoc@box
\newcommand*\pandocbounded[1]{% scales image to fit in text height/width
  \sbox\pandoc@box{#1}%
  \Gscale@div\@tempa{\textheight}{\dimexpr\ht\pandoc@box+\dp\pandoc@box\relax}%
  \Gscale@div\@tempb{\linewidth}{\wd\pandoc@box}%
  \ifdim\@tempb\p@<\@tempa\p@\let\@tempa\@tempb\fi% select the smaller of both
  \ifdim\@tempa\p@<\p@\scalebox{\@tempa}{\usebox\pandoc@box}%
  \else\usebox{\pandoc@box}%
  \fi%
}
% Set default figure placement to htbp
\def\fps@figure{htbp}
\makeatother





\setlength{\emergencystretch}{3em} % prevent overfull lines

\providecommand{\tightlist}{%
  \setlength{\itemsep}{0pt}\setlength{\parskip}{0pt}}



 


\usepackage{amsmath,amssymb,enumitem,booktabs}
\setlist{itemsep=2pt,topsep=3pt}
\usepackage{fancyhdr}\pagestyle{fancy}\fancyhf{}
\fancyhead[L]{\small DS701 Fall 2026}\fancyhead[R]{\small Homework 7}
\fancyfoot[C]{\small\thepage}
\renewcommand{\headrulewidth}{0.4pt}
\usepackage[breakable]{tcolorbox}
% --- answer-space macros (problems PDF only; no-ops in the solutions PDF) ---
% \answerspace{2in}  vertical room for handwritten work, sized from the solution
% \blank             a short fill-in rule; \blank[1.5in] for a wider one
% \answerpage        start the next problem on a fresh page (Gradescope page matching)
\newcommand{\answerspace}[1]{\par\vspace*{#1}}  % \vspace* so room is not dropped at a page break
\newcommand{\blank}[1][1.2in]{\rule{#1}{0.4pt}}
\newcommand{\answerpage}{\clearpage}
\newtcolorbox{solutionbox}{breakable,colback=blue!4,colframe=blue!45!black,title=Solution,fonttitle=\bfseries,left=4pt,right=4pt,top=3pt,bottom=3pt}
\makeatletter
\@ifpackageloaded{caption}{}{\usepackage{caption}}
\AtBeginDocument{%
\ifdefined\contentsname
  \renewcommand*\contentsname{Table of contents}
\else
  \newcommand\contentsname{Table of contents}
\fi
\ifdefined\listfigurename
  \renewcommand*\listfigurename{List of Figures}
\else
  \newcommand\listfigurename{List of Figures}
\fi
\ifdefined\listtablename
  \renewcommand*\listtablename{List of Tables}
\else
  \newcommand\listtablename{List of Tables}
\fi
\ifdefined\figurename
  \renewcommand*\figurename{Figure}
\else
  \newcommand\figurename{Figure}
\fi
\ifdefined\tablename
  \renewcommand*\tablename{Table}
\else
  \newcommand\tablename{Table}
\fi
}
\@ifpackageloaded{float}{}{\usepackage{float}}
\floatstyle{ruled}
\@ifundefined{c@chapter}{\newfloat{codelisting}{h}{lop}}{\newfloat{codelisting}{h}{lop}[chapter]}
\floatname{codelisting}{Listing}
\newcommand*\listoflistings{\listof{codelisting}{List of Listings}}
\makeatother
\makeatletter
\makeatother
\makeatletter
\@ifpackageloaded{caption}{}{\usepackage{caption}}
\@ifpackageloaded{subcaption}{}{\usepackage{subcaption}}
\makeatother
\usepackage{bookmark}
\IfFileExists{xurl.sty}{\usepackage{xurl}}{} % add URL line breaks if available
\urlstyle{same}
% fallback for those not using the hyperref driver hyperxmp:
\makeatletter
\@ifundefined{xmpquote}{\newcommand{\xmpquote}[1]{#1}}{}
\makeatother
\hypersetup{
  pdftitle={DS701 Homework 7: Least Squares{,} Logistic Loss and Regularization},
  colorlinks=true,
  linkcolor={blue},
  filecolor={Maroon},
  citecolor={Blue},
  urlcolor={Blue},
  pdfcreator={LaTeX via pandoc}}


\title{DS701 Homework 7: Least Squares, Logistic Loss and
Regularization}
\usepackage{etoolbox}
\makeatletter
\providecommand{\subtitle}[1]{% add subtitle to \maketitle
  \apptocmd{\@title}{\par {\large #1 \par}}{}{}
}
\makeatother
\subtitle{Released Mon Oct 19 · due Mon Oct 26, 11:59 pm on Gradescope}
\author{}
\date{}
\begin{document}
\maketitle


\begin{center}\LARGE\bfseries\color{blue!45!black} SOLUTIONS\end{center}
\renewcommand{\answerspace}[1]{}
\renewcommand{\answerpage}{\par\vspace{6pt}}

\answerpage

\subsection{Problem 1: Least squares by hand, and as a
projection}\label{problem-1-least-squares-by-hand-and-as-a-projection}

\emph{Practices: forming and solving the normal equations, and reading
the result geometrically.}

Fit the line \(y = \beta_0 + \beta_1 x\) to the four points
\[(x, y) \in \{(1, 2),\ (2, 3),\ (3, 5),\ (4, 6)\}.\]

\begin{enumerate}
\def\labelenumi{(\alph{enumi})}
\item
  Write the design matrix \(X\) (with a column of ones for the
  intercept) and the vector \(\mathbf{y}\). Form \(X^TX\) and
  \(X^T\mathbf{y}\), and solve the normal equations
  \(X^TX\boldsymbol\beta = X^T\mathbf{y}\) for
  \(\hat\beta_0, \hat\beta_1\). This produces a system of 2 equations in
  2 unknowns thatyou can solve directly without matrix inversion. (Both
  are exact decimals --- no rounding needed.)

  \(\hat\beta_0 =\) \blank[0.8in] \(\quad \hat\beta_1 =\) \blank[0.8in]
\end{enumerate}

\answerspace{3in}

\begin{enumerate}
\def\labelenumi{(\alph{enumi})}
\setcounter{enumi}{1}
\item
  Compute the fitted values
  \(\hat{\mathbf y} = X\hat{\boldsymbol\beta}\) and the residual vector
  \(\mathbf r = \mathbf y - \hat{\mathbf y}\). Verify
  \(X^T\mathbf r = \mathbf 0\) --- i.e.~that \(\mathbf r\) is orthogonal
  to \emph{both} columns of \(X\). Then compute TSS, RSS and \(R^2\).

  \(\hat{\mathbf y} =\) \blank[2in] \(\quad \mathbf r =\)
  \blank[1.8in]TSS \(=\) \blank[0.8in] \(\quad\) RSS \(=\)
  \blank[0.8in] \(\quad R^2 =\) \blank[0.8in]
\end{enumerate}

\answerspace{2.5in}

\begin{enumerate}
\def\labelenumi{(\alph{enumi})}
\setcounter{enumi}{2}
\tightlist
\item
  In one or two sentences: the normal equations express ``residual
  \(\perp\) column space of \(X\)''. Explain what ``\(\hat{\mathbf y}\)
  is the projection of \(\mathbf y\) onto the column space of \(X\)''
  means for these four points --- which set is \(\mathbf y\) being
  projected onto, and why the projection is the least-squares choice.
\end{enumerate}

\answerspace{2.5in}

\begin{enumerate}
\def\labelenumi{(\alph{enumi})}
\setcounter{enumi}{3}
\item
  Suppose a colleague accidentally adds the feature \(x\) a second time,
  so the design matrix becomes
  \(X' = [\mathbf 1\ \ \mathbf x\ \ \mathbf x]\) (\(4\times3\)). Show
  that \(X'^TX'\) is singular (a one-line argument or a determinant will
  do). What can you say about the least-squares coefficient vector now
  --- is it unique? --- and about the fitted values \(\hat{\mathbf y}\)?
  Give two different coefficient vectors that produce the same fit as in
  (a).

  Two coefficient vectors \((\beta_0, \beta_1, \beta_2)\):
  \blank[1.5in] and \blank[1.5in]
\end{enumerate}

\answerspace{2.5in}

\begin{solutionbox}

\begin{enumerate}
\def\labelenumi{(\alph{enumi})}
\item
  \(X = \begin{bmatrix}1&1\\1&2\\1&3\\1&4\end{bmatrix}\),
  \(\mathbf y = (2,3,5,6)^T\).
  \(X^TX = \begin{bmatrix}4&10\\10&30\end{bmatrix}\) (\(n\),
  \(\sum x_i\), \(\sum x_i^2\)),
  \(X^T\mathbf y = \begin{bmatrix}16\\47\end{bmatrix}\) (\(\sum y_i\),
  \(\sum x_iy_i = 2+6+15+24\)). Normal equations:
  \(4\beta_0 + 10\beta_1 = 16\) and \(10\beta_0 + 30\beta_1 = 47\).
  Multiply the first by \(2.5\): \(10\beta_0 + 25\beta_1 = 40\);
  subtract from the second:
  \(5\beta_1 = 7 \Rightarrow \hat\beta_1 = 1.4\), then
  \(\hat\beta_0 = (16 - 14)/4 = 0.5\). (Equivalently
  \((X^TX)^{-1} = \tfrac{1}{20}\begin{bmatrix}30&-10\\-10&4\end{bmatrix}\);
  \(\det X^TX = 120-100 = 20\).) So \(\hat y = 0.5 + 1.4x\).
\item
  \(\hat{\mathbf y} = (1.9,\,3.3,\,4.7,\,6.1)^T\);
  \(\mathbf r = (0.1,\,-0.3,\,0.3,\,-0.1)^T\).
  \(\mathbf 1^T\mathbf r = 0.1-0.3+0.3-0.1 = 0\) and
  \(\mathbf x^T\mathbf r = 0.1 - 0.6 + 0.9 - 0.4 = 0\), so
  \(X^T\mathbf r = \mathbf 0\). \(\bar y = 4\): TSS \(= 4+1+1+4 = 10\);
  RSS \(= 0.01+0.09+0.09+0.01 = 0.20\); \(R^2 = 1 - 0.20/10 = 0.98\)
  (and ESS \(= 9.8\), TSS \(=\) RSS \(+\) ESS as it must with an
  intercept).
\item
  The column space of \(X\) is the plane in \(\mathbb R^4\) spanned by
  \(\mathbf 1 = (1,1,1,1)\) and \(\mathbf x = (1,2,3,4)\) --- the set of
  all vectors of fitted values \(\beta_0\mathbf 1 + \beta_1\mathbf x\)
  any line could produce. \(\mathbf y\) does not lie in that plane (no
  line goes through all four points), so the closest achievable
  \(\hat{\mathbf y}\) is the orthogonal projection of \(\mathbf y\) onto
  it: the residual \(\mathbf y - \hat{\mathbf y}\) is perpendicular to
  the plane (that is exactly \(X^T\mathbf r = \mathbf 0\)), and
  perpendicular means shortest, i.e.~smallest
  \(\Vert\mathbf y - X\boldsymbol\beta\Vert^2\).
\item
  The second and third columns of \(X'\) are equal, so
  \(\operatorname{rank}X' = 2 < 3\) and \(X'^TX'\) (whose second and
  third rows/columns are identical) has determinant \(0\) --- singular.
  The normal equations still have solutions but not a unique one: only
  \(\beta_1 + \beta_2\) is pinned down (it must equal \(1.4\)), so there
  is a whole line of minimizers. The fitted values are unchanged ---
  \(\hat{\mathbf y}\) is the projection onto the \emph{same} column
  space, and the projection is unique even when its coordinates are not.
  Two examples: \((\beta_0,\beta_1,\beta_2) =
  (0.5, 1.4, 0)\) and \((0.5, 0.4, 1.0)\) (or \((0.5, 2.4, -1.0)\)), all
  giving \(\hat{\mathbf y} = (1.9,3.3,4.7,6.1)\). \emph{Common mistake:}
  claiming the fit itself breaks --- it is the coefficients, not the
  predictions, that become non-identifiable (and, numerically, huge and
  unstable when the columns are merely \emph{nearly} equal ---
  multicollinearity).
\end{enumerate}

\end{solutionbox}

\answerpage

\subsection{Problem 2: Logistic loss and one gradient
step}\label{problem-2-logistic-loss-and-one-gradient-step}

\emph{Practices: the sigmoid, the log-loss, and the chain rule that
makes its gradient so simple.}

Consider a one-feature logistic model
\(p(y{=}1 \mid x) = \sigma(wx + b)\) with \(\sigma(z) = 1/(1+e^{-z})\).
Current parameters: \(w = 1\), \(b = -1\). Three labelled points:
\[(x, y) \in \{(2, 1),\ (0, 0),\ (-1, 1)\}.\] The \textbf{log-loss}
(negative log-likelihood) of one example is
\(\ell_i = -\big[y_i\log p_i + (1-y_i)\log(1-p_i)\big]\) with
\(p_i = \sigma(wx_i + b)\); the total loss is \(\ell = \sum_i \ell_i\).
Minimizing \(\ell\) is the same as maximizing the likelihood \(L\). You
may use \(e^{-1}\approx0.368\), \(e^{-2}\approx0.135\),
\(\ln 0.731\approx-0.313\), \(\ln 0.119\approx-2.127\), and
\(\sigma'(z) = \sigma(z)\big(1-\sigma(z)\big)\) (take it as given, or
verify it in one line).

\begin{enumerate}
\def\labelenumi{(\alph{enumi})}
\tightlist
\item
  For each point compute \(z_i = wx_i + b\), \(p_i = \sigma(z_i)\) (3
  decimals) and \(\ell_i\). Report the total loss \(\ell\). Which point
  is the model most wrong about? You can use a calculator to compute the
  sigmoid and log-loss. Problems like this will be simplified for the
  quiz.
\end{enumerate}

{\def\LTcaptype{none} % do not increment counter
\begin{longtable}[]{@{}ccccc@{}}
\toprule\noalign{}
\(x_i\) & \(y_i\) & \(z_i\) & \(p_i\) & \(\ell_i\) \\
\midrule\noalign{}
\endhead
\bottomrule\noalign{}
\endlastfoot
\(2\) & \(1\) & \blank[0.8in] & \blank[0.8in] & \blank[0.8in] \\
\(0\) & \(0\) & \blank[0.8in] & \blank[0.8in] & \blank[0.8in] \\
\(-1\) & \(1\) & \blank[0.8in] & \blank[0.8in] & \blank[0.8in] \\
\end{longtable}
}

Total \(\ell =\) \blank[0.8in] \(\quad\) Most wrong about the point with
\(x =\) \blank[0.8in]

\answerspace{2in}
\answerpage

\begin{enumerate}
\def\labelenumi{(\alph{enumi})}
\setcounter{enumi}{1}
\tightlist
\item
  Show, for a single example, that
  \(\dfrac{\partial \ell_i}{\partial w} = (p_i - y_i)\,x_i\) and
  \(\dfrac{\partial \ell_i}{\partial b} = p_i - y_i\). (Chain rule:
  \(\ell_i \to p_i \to z_i \to w\); the given \(\sigma'\) does the heavy
  lifting.)
\end{enumerate}

\answerspace{3.5in}

\begin{enumerate}
\def\labelenumi{(\alph{enumi})}
\setcounter{enumi}{2}
\item
  Using (b), compute \(\partial\ell/\partial w\) and
  \(\partial\ell/\partial b\) summed over the three points, and take one
  gradient-\textbf{descent} step with learning rate \(\alpha = 0.5\):
  \(w \leftarrow w - \alpha\,\partial\ell/\partial w\),
  \(b \leftarrow b - \alpha\,\partial\ell/\partial b\). Report the new
  \((w, b)\) to 3 decimals. In which direction did \(b\) move, and why
  does that make sense given (a)?

  \(\partial\ell/\partial w =\)
  \blank[0.8in] \(\quad \partial\ell/\partial b =\)
  \blank[0.8in] \(\quad\) new \((w, b) =\) \blank[1.5in]
\end{enumerate}

\answerspace{2.5in}
\answerpage

\begin{enumerate}
\def\labelenumi{(\alph{enumi})}
\setcounter{enumi}{3}
\tightlist
\item
  (Optional) A classmate proposes fitting the same model by minimizing
  squared error \(\sum_i (p_i - y_i)^2\) instead. Give two reasons, in a
  short paragraph, why the lecture fits logistic regression by maximum
  likelihood (i.e.~minimizes the log-loss) rather than squared error.
  Hint: what does the least-squares criterion assume about the noise,
  and how large a penalty does each loss assign to a confident wrong
  prediction such as \(p_i = 0.99\) when \(y_i = 0\)?
\end{enumerate}

\answerspace{1.8in}

\begin{solutionbox}

\begin{enumerate}
\def\labelenumi{(\alph{enumi})}
\item
  \(x=2\): \(z = 1\), \(p = 1/(1+0.368) = 0.731\), \(y=1\):
  \(\ell_1 = -\ln 0.731 = 0.313\). \(x=0\): \(z=-1\),
  \(p = 1/(1+2.718) = 0.269\), \(y=0\):
  \(\ell_2 = -\ln(1-0.269) = -\ln 0.731 = 0.313\). \(x=-1\): \(z=-2\),
  \(p = 1/(1+7.389) = 0.119\), \(y=1\): \(\ell_3 = -\ln 0.119 = 2.127\).
  Total \(\ell = 0.313+0.313+2.127 = 2.753\). The third point: the model
  gives the label-1 point only a \(12\%\) chance, and its loss
  dominates.
\item
  With \(\ell_i = -y_i\log p_i - (1-y_i)\log(1-p_i)\):
  \(\partial\ell_i/\partial p_i = -y_i/p_i + (1-y_i)/(1-p_i)\);
  \(\partial p_i/\partial z_i = p_i(1-p_i)\);
  \(\partial z_i/\partial w = x_i\), \(\partial z_i/\partial b = 1\).
  Multiply:
  \[\frac{\partial\ell_i}{\partial z_i} = \Big(-\frac{y_i}{p_i} + \frac{1-y_i}{1-p_i}\Big)p_i(1-p_i)
  = -y_i(1-p_i) + (1-y_i)p_i = p_i - y_i,\] so
  \(\partial\ell_i/\partial w = (p_i-y_i)x_i\) and
  \(\partial\ell_i/\partial b = p_i-y_i\). (Verifying \(\sigma'\):
  \(\sigma = (1+e^{-z})^{-1}\), so
  \(\sigma' = e^{-z}(1+e^{-z})^{-2} = \sigma\cdot\frac{e^{-z}}{1+e^{-z}} = \sigma(1-\sigma)\).)
\item
  \(p_i - y_i\): \(-0.269,\ +0.269,\ -0.881\).
  \(\partial\ell/\partial w = (-0.269)(2) + (0.269)(0) + (-0.881)(-1) = -0.538 + 0.881 = 0.343\);
  \(\partial\ell/\partial b = -0.269 + 0.269 - 0.881 = -0.881\). Step:
  \(w = 1 - 0.5(0.343) = 0.829\) (0.828 with unrounded intermediates ---
  either accepted), \(b = -1 - 0.5(-0.881) = -0.560\). \(b\) went
  \emph{up}: two of the three points are label 1 and the model is
  under-predicting them (net \(\sum(p_i-y_i) < 0\)), so raising the base
  rate lowers the loss (it drops from \(2.75\) to about \(2.35\)).
  \emph{Common mistake:} adding \(\alpha\nabla\ell\) --- that is ascent
  on the \emph{loss}. Written as gradient ascent, the update adds
  \(\alpha\) times the gradient of the \emph{log-likelihood} --- the
  same move.
\item
  \begin{enumerate}
  \def\labelenumii{(\roman{enumii})}
  \tightlist
  \item
    Least squares is maximum likelihood only under Gaussian noise; a 0/1
    label is a Bernoulli outcome, and its likelihood is
    \(\prod p_i^{y_i}(1-p_i)^{1-y_i}\) --- maximizing that is exactly
    minimizing the log-loss, so log-loss is the principled objective and
    its minimizer gives calibrated probabilities. (ii) Squared error
    caps the penalty for a confident wrong answer at 1 (\(p=0.99, y=0\):
    \(0.98\)), while log-loss charges \(-\ln 0.01 = 4.6\) and grows
    without bound --- it refuses to let the model be confidently wrong.
    A third acceptable reason: with the sigmoid inside, squared error is
    non-convex in \((w,b)\) with flat regions where the gradient
    \(\propto p(1-p)\) vanishes for saturated wrong predictions, whereas
    the log-loss gradient \((p-y)x\) stays large exactly when the error
    is large --- gradient descent behaves well.
  \end{enumerate}
\end{enumerate}

\end{solutionbox}

\answerpage

\subsection{Problem 3: Ridge shrinks, LASSO
zeroes}\label{problem-3-ridge-shrinks-lasso-zeroes}

\emph{Practices: what the L2 and L1 penalties actually do to a
coefficient, and how to choose the penalty weight.}

Throughout, \(\lambda \ge 0\) is the penalty weight.

\begin{enumerate}
\def\labelenumi{(\alph{enumi})}
\item
  \textbf{Ridge in one dimension.} For a no-intercept model
  \(y \approx \beta x\) with data
  \((x, y) \in \{(1,1),\ (2,4),\ (3,7)\}\), the ridge objective is
  \(J(\beta) = \sum_i (y_i - \beta x_i)^2 + \lambda\beta^2\). By setting
  \(dJ/d\beta = 0\), show that
  \(\hat\beta_\lambda = \dfrac{\sum_i x_iy_i}{\sum_i x_i^2 + \lambda}\)
  (this is \((X^TX + \lambda I)^{-1}X^T\mathbf y\) with \(X\) a single
  column). Evaluate it for \(\lambda = 0, 1, 10\) (3 decimals) and
  describe in a sentence what happens as \(\lambda\) grows. Can any
  finite \(\lambda\) make \(\hat\beta_\lambda\) exactly \(0\)?

  \(\hat\beta_0 =\) \blank[0.8in] \(\quad \hat\beta_1 =\)
  \blank[0.8in] \(\quad \hat\beta_{10} =\) \blank[0.8in] \(\quad\)
  exactly \(0\)? \blank[0.6in]
\end{enumerate}

\answerspace{2.5in}
\newpage

\begin{enumerate}
\def\labelenumi{(\alph{enumi})}
\setcounter{enumi}{1}
\item
  \textbf{LASSO in one dimension.} The building block of the LASSO is
  the scalar problem
  \[\min_\beta\ \tfrac12(\beta - z)^2 + \lambda|\beta|,\] where \(z\) is
  the unpenalized (least-squares) value. Assume first \(z > 0\). Show
  that the minimizer is \(\hat\beta = z - \lambda\) if \(z > \lambda\)
  and \(\hat\beta = 0\) if \(z \le \lambda\). (Argue in two cases: if
  the minimizer is \(\beta > 0\), differentiate; if \(\beta<0\), show it
  cannot be a minimizer when \(z>0\); the remaining candidate is
  \(\beta = 0\).) By symmetry the general answer is
  \(\hat\beta = \operatorname{sign}(z)\max(|z|-\lambda, 0)\) --- the
  \emph{soft threshold}. Evaluate it for \(\lambda = 1\) and \(z = 3\),
  \(z = 0.5\), \(z = -2\).

  LASSO, \(\lambda=1\): \(z=3 \to\) \blank[0.8in] \(\quad z=0.5 \to\)
  \blank[0.8in] \(\quad z=-2 \to\) \blank[0.8in]
\end{enumerate}

\answerspace{3in}

\begin{enumerate}
\def\labelenumi{(\alph{enumi})}
\setcounter{enumi}{2}
\item
  The ridge version of the same scalar problem,
  \(\tfrac12(\beta - z)^2 + \lambda\beta^2\), has minimizer
  \(\hat\beta = z/(1+2\lambda)\) (you need not derive it). Evaluate it
  for \(\lambda = 1\) and the same three \(z\), and then explain in two
  or three sentences, using what you found in (a)--(c), why the L1
  penalty produces coefficients that are \emph{exactly} zero while the
  L2 penalty only shrinks them. (Think about the slope of \(|\beta|\)
  versus \(\beta^2\) at \(\beta = 0\).)

  Ridge, \(\lambda=1\): \(z=3 \to\) \blank[0.8in] \(\quad z=0.5 \to\)
  \blank[0.8in] \(\quad z=-2 \to\) \blank[0.8in]
\end{enumerate}

\answerspace{2in}
\answerpage

\begin{enumerate}
\def\labelenumi{(\alph{enumi})}
\setcounter{enumi}{3}
\tightlist
\item
  \textbf{Scenario.} An analyst fits LASSO to predict house price from
  \texttt{sqft} (values in the thousands), \texttt{bedrooms} (1--5) and
  \texttt{age\_years} (0--100), without standardizing, and chooses
  \(\lambda\) as the value that minimizes the \emph{training} RSS. State
  the two things wrong with this procedure and what to do instead (one
  or two sentences each).
\end{enumerate}

\answerspace{2in}

\begin{solutionbox}

\begin{enumerate}
\def\labelenumi{(\alph{enumi})}
\item
  \(dJ/d\beta = -2\sum_i x_i(y_i - \beta x_i) + 2\lambda\beta = 0
  \Rightarrow \beta(\sum x_i^2 + \lambda) = \sum x_iy_i\). Here
  \(\sum x_iy_i = 1 + 8 + 21 = 30\) and \(\sum x_i^2 = 14\), so
  \(\hat\beta_0 = 30/14 = 2.143\), \(\hat\beta_1 = 30/15 = 2.000\),
  \(\hat\beta_{10} = 30/24 = 1.250\). The estimate is pulled
  monotonically toward \(0\) as \(\lambda\) grows (the denominator grows
  while the numerator is fixed), but \(30/(14+\lambda) > 0\) for every
  finite \(\lambda\): ridge shrinks, never zeroes.
\item
  Let \(f(\beta) = \tfrac12(\beta-z)^2 + \lambda|\beta|\), \(z > 0\).
  \emph{Case \(\beta>0\):} \(f = \tfrac12(\beta-z)^2 + \lambda\beta\),
  \(f' = \beta - z + \lambda = 0
  \Rightarrow \beta = z - \lambda\), valid (positive) only if
  \(z > \lambda\); \(f''=1>0\) so it is a minimum. \emph{Case
  \(\beta<0\):} \(f' = \beta - z - \lambda < 0\) throughout, so \(f\) is
  decreasing on \((-\infty,0)\) and no negative point is a minimizer ---
  moving toward \(0\) always helps. \emph{Case \(\beta = 0\):} the only
  remaining candidate; when \(z \le \lambda\) the right-hand slope at
  \(0\) is \(-z + \lambda \ge 0\) and the left-hand slope is
  \(-z-\lambda < 0\), so \(f\) increases as you move away from \(0\) in
  either direction --- \(0\) is the minimizer. Hence
  \(\hat\beta = z-\lambda\) if \(z>\lambda\), else \(0\); with the sign
  case, \(\hat\beta = \operatorname{sign}(z)\max(|z|-\lambda,0)\).
  \(\lambda=1\): \(z=3 \to 2\); \(z=0.5 \to 0\); \(z=-2 \to -1\).
\item
  Ridge, \(\lambda = 1\): \(z=3\to1\); \(z=0.5\to0.167\);
  \(z=-2\to-0.667\) --- every value is scaled by \(1/3\), none is zero.
  Reason: near \(\beta=0\) the L2 penalty \(\lambda\beta^2\) is flat
  (derivative \(2\lambda\beta \to 0\)), so it exerts no force at zero
  and any nonzero \(z\) pulls the coefficient off zero --- the L2
  optimum balances two derivatives and lands at \(z/(1+2\lambda)\),
  small but nonzero. The L1 penalty \(\lambda|\beta|\) has a kink at
  \(0\) with slope \(\pm\lambda\) on either side, so as long as the pull
  of the data \(|z|\) is below \(\lambda\), the penalty wins and
  \(\hat\beta = 0\) exactly. That is why LASSO does variable selection
  and ridge does not.
\item
  \begin{enumerate}
  \def\labelenumii{(\arabic{enumii})}
  \tightlist
  \item
    \emph{No standardization:} the penalty
    \(\lambda\Vert\boldsymbol\beta\Vert_1\) charges every coefficient
    the same price per unit, but a coefficient on \texttt{sqft}
    (thousands) is naturally small and one on \texttt{bedrooms} (1--5)
    naturally large, so LASSO would zero out \texttt{sqft} first for
    reasons of units, not relevance. Fix: center \(\mathbf y\) and
    standardize each column of \(X\) (zero mean, unit variance) before
    penalizing. (2) \emph{Choosing \(\lambda\) by training RSS:} the
    training RSS is minimized at \(\lambda = 0\) (any penalty can only
    raise it), so the rule always picks plain least squares and defeats
    the purpose. Fix: choose \(\lambda\) by cross-validation / held-out
    error, which is what regularization is trying to improve.
  \end{enumerate}
\end{enumerate}

\end{solutionbox}




\end{document}
